% 保存后由 LaTeX Workshop 使用 XeLaTeX 编译，也支持 LuaLaTeX。
% 终端编译：latexmk -xelatex -interaction=nonstopmode -halt-on-error -outdir=PDF example.tex
\documentclass[UTF8,fontset=fandol,a4paper,11pt]{ctexart}
\usepackage[margin=2.5cm]{geometry}
\usepackage{amsmath,amssymb}
\usepackage{booktabs}
\usepackage{siunitx}
\usepackage[unicode,hidelinks]{hyperref}

\title{WebLaTeX CN 中文科技写作示例}
\author{你的名字}
\date{\today}

\begin{document}
\maketitle

\begin{abstract}
这是一份可以独立编译的中文示例，展示中文排版、数学公式、科学单位、表格与参考文献。
修改本文件并保存，即可在 \texttt{PDF/example.pdf} 中查看结果。
\end{abstract}

\noindent\textbf{关键词：} 中文排版；LaTeX；科技写作

\section{开始写作}
本示例使用随环境提供的 Fandol 字体，无需安装额外字体。
可以直接输入中文与 English，也可以使用\textbf{粗体}和\emph{强调}。
LaTeX 的基本使用可参考文献~\cite{knuth1984}。

\section{数学公式与单位}
行内公式可以写成 $E=mc^2$。带编号的公式如下：
\begin{equation}
  \int_0^1 x^2\,\mathrm{d}x = \frac{1}{3}.
  \label{eq:integral}
\end{equation}
公式~\eqref{eq:integral} 可通过标签交叉引用，无需手动填写编号。
多行推导可以使用 \texttt{align} 环境：
\begin{align}
  f(x) &= x^2 + 2x + 1, \\
  f'(x) &= 2x + 2.
\end{align}
科学单位使用 \texttt{siunitx}，例如温度 \qty{300}{\kelvin}、
长度 \qty{12.5}{\milli\metre} 和速度 \qty{3.0}{\metre\per\second}。

\section{表格}
表~\ref{tab:demo} 中的数据仅用于演示排版，不代表真实实验结果。
\begin{table}[htbp]
  \centering
  \caption{排版演示数据}
  \label{tab:demo}
  \begin{tabular}{lrr}
    \toprule
    条目 & 数量 & 得分 \\
    \midrule
    示例 A & 10 & 0.85 \\
    示例 B & 20 & 0.92 \\
    \bottomrule
  \end{tabular}
\end{table}

\section{下一步}
将标题、作者和正文替换为自己的内容即可开始写作。
本文件内置参考文献，因此无需另建 \texttt{.bib} 文件。
正式写作时，也可以改用环境中已安装的 BibTeX 或 Biber 管理文献。

\begin{thebibliography}{9}
  \bibitem{knuth1984}
  Donald E. Knuth.
  \emph{The TeXbook}.
  Addison-Wesley, 1984.
\end{thebibliography}

\end{document}
